\documentclass[11pt]{article}

\usepackage[a4paper,margin=2.6cm]{geometry}
\usepackage[T1]{fontenc}
\usepackage[utf8]{inputenc}
\usepackage{lmodern}
\usepackage{microtype}
\usepackage{amsmath,amssymb,mathtools}
\usepackage{booktabs}
\usepackage{array}
\usepackage{enumitem}
\usepackage{graphicx}
\usepackage{tikz}
\usetikzlibrary{arrows.meta,positioning,fit}
\usepackage[backend=biber,style=authoryear,natbib=true,maxcitenames=2,maxbibnames=99]{biblatex}
\addbibresource{references.bib}
\usepackage[hidelinks]{hyperref}
\usepackage[nameinlink,noabbrev]{cleveref}
\emergencystretch=2em

\newcommand{\Prob}{\mathbb{P}}
\newcommand{\E}{\mathbb{E}}
\newcommand{\R}{\mathbb{R}}
\newcommand{\calO}{\mathcal{O}}
\newcommand{\calS}{\mathcal{S}}
\newcommand{\calC}{\mathcal{C}}
\newcommand{\method}{\textsc{FAMS}}

\title{Failure-Aware Multimodal Behavioural Sensing:\\
Probabilistic State Inference, Person Attribution, and Sensor Reduction in Smart Homes}
\author{Diogo Ribeiro\\
\small \textit{ESMAD, Instituto Politécnico do Porto}}
\date{}

\begin{document}
\maketitle

\begin{abstract}
Ambient assisted-living systems infer behaviour from sparse, heterogeneous, and imperfect sensor streams, yet sensor activation is often treated too directly as behavioural truth. We present a failure-aware multimodal framework in which sensor events are uncertain observations of latent behavioural state. The framework combines hardware-neutral observation semantics, sensor-health reliability, occupancy-aware resident attribution, continuous-time latent-state filtering, adaptive baselines, restrained alerting, and paired sensor ablation. We evaluated five pre-specified hypotheses in a frozen confirmatory simulation with $N=200$ paired household trajectories; the run passed its Monte Carlo precision and provenance gates. Random missingness produced graded loss of balanced accuracy, although abstention did not increase materially. The frozen five-sensor deployment failed non-inferiority against the ten-sensor deployment by a large margin: the balanced-accuracy gap was $0.1548$ (95\% CI $[0.1531,0.1564]$) against a $0.02$ margin. A secondary eight-sensor configuration was much closer to the full deployment, with a gap of $0.00529$. Occupancy-aware attribution helped under several contamination regimes, particularly carer-related scenarios, but harmed balanced accuracy when wearable and beacon evidence were unavailable. Failure-aware reliability weighting improved log loss while worsening balanced accuracy, Brier score, and calibration error, demonstrating a genuine scoring trade-off rather than uniform dominance. All four pre-specified sensor-interaction terms were positive. These results are simulator-derived and are not estimates of field performance. Development-stage evidence from 22 real annotated homes, summarised here rather than omitted, indicates that the simulator represents an easier problem than the instrumentation supports --- median balanced accuracy $0.420$ against a $0.607$ supervised ceiling on the same sensor information --- and independently corroborates the H1 abstention finding. The contribution is therefore methodological and diagnostic: fewer sensors can preserve information, but aggressive reduction and reliability weighting have failure modes that must remain visible rather than being hidden by post-hoc model selection.
\end{abstract}

\noindent\textbf{Keywords:} ambient assisted living; behavioural sensing; sensor fusion; uncertainty; smart homes; missing data; multi-resident activity recognition; sensor ablation

\section{Introduction}

Ambient assisted living (AAL) offers a route to continuous, low-burden monitoring in the home, but the useful signal is indirect. Motion detectors, contact switches, bed sensors, wearables, and radar systems measure physical phenomena; they do not directly observe activities of daily living, health states, or clinical deterioration. Recent reviews describe AAL design as a system-level trade-off between sensing quality, privacy, usability, reliability, communication, cost, and decision-making \citep{zieni2025aal}. Statistical work on passive sensing has shown that interpretable models can capture within-home temporal structure and cross-sensor dependence \citep{gillam2022athome}, while smart-home activity recognition has developed increasingly sophisticated sequential and multi-person models \citep{asghari2019hhmm,wang2022multiperson}.

Four methodological problems motivate the present framework. First, sensor semantics are often stronger than the evidence: a refrigerator-door activation is evidence of interaction, not proof of food intake. Second, missingness and behavioural quiet are confounded unless sensor-health semantics are represented explicitly. Third, activity in the home is not necessarily activity by the monitored resident. Fourth, more sensors are not automatically better because installation burden, maintenance, privacy, and failure surfaces also increase.

We therefore model a posterior distribution over latent behavioural states,
\begin{equation}
\Prob(Z_t=z\mid\calO_{1:t}),
\end{equation}
with sensor quality, health reliability, and resident attribution tempering the evidential contribution of observations. A failed sensor contributes no state preference; it does not vote for inactivity.

The paper makes five methodological contributions: a canonical observation model separating measurement from interpretation; reliability-aware multimodal filtering; occupancy-aware resident attribution; adaptive personal baselines for behavioural change; and paired sensor-ablation and interaction analysis. The second is the one this manuscript is least able to support: abstention is specified and implemented, but every measurement of whether it works is negative, in the confirmatory simulation, on the real development recordings, and in the external validation. We report it as a contribution to the architecture and as a negative result about the formulation. Unlike the earlier development manuscript, the present version reports the frozen confirmatory simulator study. Its confirmatory conclusions are simulator-derived, but this manuscript is no longer simulator-only: it also reports development evidence from 22 real annotated homes and the outcome of the frozen external validation, at the evidential levels set out below.

\paragraph{Evidence status.}
Three levels of evidence appear here and are not interchangeable. The confirmatory results are simulator-derived, frozen before execution, and estimate nothing about field performance. The real-recording results are development evidence: those homes were inspected repeatedly while diagnosing failure modes, so they are not a held-out test. The external-validation result is confirmatory on real data, scored once against a cohort frozen before any outcome was seen, and is the only figure here that carries that status.

\section{Related work}

Passive and ambient sensing is attractive because it can observe routine without repeated user interaction. Transparent household-specific probabilistic models can capture temporal and cross-sensor structure \citep{gillam2022athome}. Sequential latent-state models are natural for activity recognition because activities exhibit persistence and transitions, while change-point methods offer a complementary view focused on transitions \citep{asghari2019hhmm,bermejo2021cpd,culman2020cpam}.

Person attribution is a distinct inference problem from activity recognition. ARAS provides multi-resident smart-home data with labelled activities \citep{alemdar2013aras}, and multi-person activity-recognition work shows that associating observations with residents can materially change recognition performance \citep{wang2022multiperson}. Our scope is narrower: one monitored resident is modelled explicitly, while other occupants are contextual contamination rather than targets of biometric identification.

Privacy-preserving richer sensing, including millimetre-wave radar, is increasingly relevant because it can provide presence, motion, posture, and tracking information without optical imagery \citep{chen2026mmwave,samimi2026fmcw}. The relevant design question is not simply whether such a modality improves accuracy, but whether it can replace narrower sensors without obscuring uncertainty or creating brittle dependence on one high-information stream.

\section{Methodological framework}
\label{sec:method}

A canonical observation is represented as
\begin{equation}
O_i=(t_i,s_i,m_i,k_i,x_i,u_i,q_i,c_i,\rho_i),
\end{equation}
where the fields encode time, sensor identity, modality, temporal semantics, value, unit, device quality, confidence, and provenance/context. Events, persistent states, and sampled quantities are kept distinct so that absent observations are not silently converted into measured zeros.

For sensor $s$ at time $t$, the health layer estimates
\begin{equation}
r_s(t)\in[0,1].
\end{equation}
Reliability depends on declared reporting semantics. A zero reliability weight means no state preference, not evidence for inactivity.

Occupancy context is represented by $C_t$, and ambient evidence receives an attribution factor
\begin{equation}
a_s(t)=\Prob(\text{resident generated evidence from }s\mid\calO_{1:t}).
\end{equation}
Person-bound sensors can remain fully attributable when valid, whereas ambient evidence is discounted when another occupant may be its source.

Between observations, the latent-state belief evolves under a continuous-time Markov generator $Q$,
\begin{equation}
\boldsymbol p(t+\Delta)=\boldsymbol p(t)\exp(Q\Delta).
\end{equation}
For observations $B_t$, the tempered update is
\begin{equation}
p_t(z)\propto p_{t^-}(z)\prod_{s\in B_t}L_s(z)^{\omega_s(t)},
\end{equation}
with
\begin{equation}
\omega_s(t)=\lambda_s q_s(t)c_s(t)r_s(t)a_s(t).
\end{equation}
This power-likelihood form reduces the influence of uncertain evidence without pretending that uncertainty is a deterministic negative observation, and is related to general Bayesian updating \citep{bissiri2016general}.

Daily behavioural summaries integrate posterior state occupancy rather than only argmax labels. Poorly observed days can be excluded from baseline learning, and robust deviation scores and change-point methods separate ordinary variability from persistent shifts \citep{killick2012pelt}. Alerting remains downstream of state and change inference so that statistical novelty is not automatically equated with a delivered alert.

For a deployment sensor set $S$ and metric $M(S)$, sensor marginal effects and interactions are defined by
\begin{equation}
\Delta_j=M(S)-M(S\setminus\{j\}),
\end{equation}
and
\begin{equation}
I_{ij}=\big[M(S)-M(S\setminus\{i,j\})\big]
-\big[M(S)-M(S\setminus\{i\})\big]
-\big[M(S)-M(S\setminus\{j\})\big].
\end{equation}
Positive $I_{ij}$ indicates complementarity under this definition.

\section{Experimental design}
\label{sec:experiments}

\subsection{Metrics}
\label{sec:metrics}

Balanced accuracy is the mean per-class recall over the states present in the
truth, which prevents a model that always predicts the majority state from
scoring well; macro F1 is the unweighted mean F1 over the same classes.
Multiclass log loss and the multiclass Brier score are computed from the full
posterior rather than from the argmax label. Expected calibration error uses ten
equal-width confidence bins,
\begin{equation}
\mathrm{ECE}=\sum_{b=1}^{10}\frac{n_b}{n}\big|\mathrm{acc}(b)-\mathrm{conf}(b)\big|,
\end{equation}
where $n_b$ is the number of estimates whose leading-state confidence falls in
bin $b$, and $\mathrm{acc}(b)$ and $\mathrm{conf}(b)$ are the mean correctness
and mean confidence within it.

The abstention convention matters for reading H1 and is applied consistently: a
reported \textsc{unknown} is never scored correct, because \textsc{unknown} is
never a true label. Abstention therefore costs recall, and abstention rate is
reported alongside the accuracy metrics so that a model which declines usefully
can be distinguished from one that is simply wrong.

\subsection{Simulator and frozen design}

The simulator generates longitudinal household trajectories with known latent state, resident location, sleep, outings, visitors, carer visits, and sensor observations. The generative process is deliberately structurally different from the estimator. Sensor degradation is injected separately from behavioural generation, allowing paired counterfactual comparisons on the same household trajectory.

The confirmatory study was frozen before results were examined. The independent replication unit is one simulated household trajectory, with initial production size $N=200$ and seeds
\begin{equation}
s_i=10000+4i,\qquad i=0,\ldots,N-1.
\end{equation}
Replication could increase only for the pre-specified Monte Carlo precision criterion, up to $N=1000$, without inspecting hypothesis direction or changing the design. The target was maximum H2 balanced-accuracy MCSE $\leq0.002$ across the pre-specified reduced configurations. The observed maximum MCSE was $0.001125$, so the experiment stopped at the initial $N=200$.

The five frozen hypotheses, their contrasts, and their pre-specified criteria
are given in \cref{tab:hypotheses}. All uncertainty calculations are over
households, never time points. A negative hypothesis result remains
scientifically valid if the frozen provenance and precision gates pass.

\begin{table}[t]
\centering
\small
\caption{Frozen confirmatory hypotheses. Only H2 carried a numeric decision
gate; the remaining hypotheses were pre-specified as directional, with the
secondary metrics supplying uncertainty and safety context rather than
additional gates.}
\label{tab:hypotheses}
\begin{tabular}{@{}lp{3.5cm}p{3.6cm}p{3.1cm}@{}}
\toprule
& Frozen contrast & Primary metric & Pre-specified criterion \\
\midrule
H1 & Random record loss at $0$, $5$, $10$, $20$, $40\%$ & Balanced accuracy, with log loss, Brier, ECE, abstention rate & Graded rather than cliff-like loss \\
H2 & Five-sensor versus ten-sensor deployment, paired by seed & $D_{H2}=\mathrm{BA}_{\mathrm{full}}-\mathrm{BA}_{\mathrm{reduced}}$ & $U_{0.95}(D_{H2})<0.02$ \\
H3 & Occupancy-aware versus naive attribution, eight contamination scenarios & Balanced accuracy and probabilistic scores & Scenario-specific benefit \\
H4 & Failure-aware versus health-naive weighting on identical degraded records & Balanced accuracy, log loss, Brier, ECE & Direction of paired difference \\
H5 & Four pre-specified sensor pairs & Interaction $I_{ij}$ & Non-zero \\
\bottomrule
\end{tabular}
\end{table}

\paragraph{Frozen design parameters.}
Each replication simulates 70 days of a single household on a 15-minute
inference grid. The latent ontology has seven states: away, home-active,
home-inactive, sleeping, bed-awake, bathroom-activity, and kitchen-activity. The
full deployment comprises ten sensors: \texttt{front\_door}; four motion sensors
in bedroom, bathroom, kitchen, and living room; \texttt{fridge\_contact};
\texttt{bed\_pressure}; \texttt{living\_radar}; \texttt{wearable\_motion}; and
\texttt{resident\_beacon}. The frozen H2 reduced configuration retains
\texttt{front\_door}, \texttt{living\_radar}, \texttt{bed\_pressure},
\texttt{wearable\_motion}, and \texttt{resident\_beacon}; the secondary
eight-sensor configuration is the six object sensors plus
\texttt{wearable\_motion} and \texttt{resident\_beacon}. The H4 degraded regime
applies 30\% random record loss together with a 14-day \texttt{bed\_pressure}
outage beginning on day 35, and the health-naive control sees exactly the same
observations with all reliability weights forced to one, so the contrast
isolates the use of health reliability rather than the missingness itself. All
intervals are 95\% household-paired bootstrap intervals over 5000 resamples.

\input{confirmatory_results.tex}

\section{Discussion}

\subsection{Graceful degradation does not imply appropriate abstention}
\label{sec:h1abstention}

The H1 result separates two ideas that are often conflated. Balanced accuracy degrades gradually as random records are removed, which is preferable to a cliff-like failure. However, the inference system does not respond to this loss by abstaining materially more often, and the probabilistic metrics are not monotone in the intuitive direction. In particular, log loss improves while discrimination worsens. The architectural rule that missing sensors should not vote for inactivity is therefore necessary but not sufficient for an uncertainty mechanism that degrades in the way a downstream decision-maker might expect. Coverage-aware abstention remains a specific target for external validation and future model revision.

\subsection{Sensor reduction is a frontier, not a five-sensor success}

The primary H2 result is unambiguously negative. The frozen five-sensor deployment loses roughly 0.155 balanced-accuracy points relative to the full ten-sensor deployment, far outside the $0.02$ non-inferiority margin. The study therefore does not support the claim that the selected five-sensor kit preserves state-inference performance.

At the same time, the secondary eight-sensor result shows that this negative result should not be simplified into ``more sensors are always better''. The eight-sensor configuration lies close to the full system, while six-, five-, three-, and two-sensor configurations lose substantially more information. The scientifically useful object is consequently a sensor-information frontier. The secondary configuration is evidence about that frontier, not a substitute primary hypothesis chosen after seeing the result.

\subsection{Attribution is useful only when attribution evidence is itself adequate}

H3 supports the idea that occupancy-aware attribution can matter under contamination, especially in carer-related regimes. Yet the reversal when wearable and resident-beacon evidence are removed is equally important. Attribution is a probabilistic nuisance-variable correction, not a universally beneficial preprocessing step. If the evidence identifying the resident is weak, discounting ambient evidence can remove useful signal from the wrong person hypothesis. External validation should therefore report attribution benefit jointly with the quality of the evidence used to perform attribution.

\subsection{Failure-aware weighting changes the error distribution rather than dominating}

H4 is a direct warning against evaluating reliability weighting with a single score. Failure-aware weighting improves log loss but worsens balanced accuracy, Brier score, and ECE on the frozen degraded-sensing regime. This means the weighting reduces some severe probability errors while producing worse average classification and calibration behaviour. The correct interpretation is a scoring trade-off, not a general robustness victory. Any operational deployment would need a loss function tied to the downstream decision rather than assuming that lower confidence under failure is automatically better in every statistical sense.

\subsection{Sensor interactions are empirically non-additive in the simulator}

All four H5 interaction terms are positive. The strongest interaction involves wearable motion and the resident beacon, with a smaller but clear interaction between bed pressure and wearable motion. These results are consistent with the idea that identity-bearing and ambient modalities can complement one another. They also explain why raw sensor count is a poor proxy for information content: what matters is not only which measurements are present but how their evidential roles combine.

\section{Limitations}
\label{sec:limitations}

The dominant limitation is external validity. The simulator was written by the same project as the estimator. Structural separation, guard tests, and distinct generative assumptions reduce circularity but cannot remove model-world bias. Simulator seeds vary stochastic realisations under one assumed world; they do not represent the structural heterogeneity of real households.

Second, several inference parameters are declared rather than estimated from independent field data, including state dwell assumptions, occupancy priors, evidence weights, and alert thresholds. Parameter uncertainty is not propagated. Third, occupancy and behavioural-state layers share evidence, so their errors are correlated even though the implementation treats the stages sequentially. Fourth, the ontology represents one monitored resident rather than a complete multi-resident latent-state system. Fifth, late observations beyond the configured tolerance are dropped rather than retrospectively smoothing past states.

Fifth, and with the widest consequences for the claims above, the abstention mechanism does not function. It barely fired under 40\% missingness here, its stated confidence carried almost no information about correctness on real recordings, and it did not move in the external validation. No threshold setting repairs a confidence signal that is only weakly related to correctness, so the framework's claim to know when its evidence is insufficient is contradicted by every body of data we have.

Sixth, the simulator does not merely flatter the framework; it represents an easier problem. Its balanced accuracy exceeds the ceiling a supervised classifier reaches on real instrumentation given the same sensors, so simulator figures cannot be read as optimistic estimates of the same quantity.

Finally, the confirmatory intervals are precise conditional on the simulator and paired household design; precision does not imply transportability. The narrow H2 interval establishes that the five-sensor gap is not Monte Carlo noise under this model world. It does not establish the size of that gap in real homes.

\section{Development evidence from real recordings}
\label{sec:realdev}

This manuscript reports the frozen confirmatory simulation, and the results
above are simulator-derived by construction. Development work on real
recordings has since been carried out and is reported in full in the companion
development manuscript. We summarise it here rather than omit it, because two of
its findings bear directly on how the confirmatory results should be read, and
because a confirmatory paper reporting only the simulator would present a more
favourable picture than this project's own evidence supports.

The evidence is development-grade and cannot be promoted to confirmatory: the 22
annotated single-resident homes involved were inspected repeatedly while
diagnosing failure modes, so they are not a held-out test. Nothing in this
section revises a frozen confirmatory outcome, and no result reported here was
used to select the confirmatory design, which was frozen first.

Two findings matter for this manuscript. First, median balanced accuracy on
those recordings was $0.420$, against $0.816$ in simulation, while a supervised
classifier given the same sensor information reached only $0.607$. The simulator
therefore represents an easier problem than the instrumentation supports for any
method, rather than an optimistic estimate of the same quantity. The
transportability caveat in \cref{sec:limitations} is consequently stronger than a
statement about interval width: the gap lies in the problem, not only in the
precision.

Second, the real recordings corroborate the H1 abstention finding of
\cref{sec:h1abstention} by an independent route. In simulation the mechanism
barely fired at all, with mean abstention rates of $9.7\times10^{-6}$ under no
missingness and $5.3\times10^{-5}$ at 40\%. On real recordings it did fire, on
$2.2\%$ of steps, but its confidence carried almost no information about
correctness: correct and incorrect answers were separated by $0.073$, and the
relationship inverted above $0.95$, so no threshold setting repairs it. The two
failures differ in character --- near-silence in simulation, uninformative
confidence on real data --- but they agree on the conclusion, and it is the
conclusion this framework can least afford: the system does not know when its
evidence is insufficient.

\section{External validation}
\label{sec:external}

External validation should proceed without redesigning the model around the simulator results. CASAS provides annotated smart-home datasets spanning single-resident and multi-resident settings \citep{cook2013casas}, while ARAS provides month-long labelled recordings from two multi-resident homes \citep{alemdar2013aras}. These datasets can test state recognition, contamination, temporal generalisation, and robustness to real sensor sparsity.

The external-validation workflow should freeze the ontology, observation semantics, metrics, and principal hypotheses; implement dataset adapters without modifying downstream inference; fit only parameters explicitly designated as trainable using household-separated training data; evaluate on held-out homes or held-out periods rather than random time points; and report discrimination, calibration, abstention, attribution evidence quality, and failure/coverage strata. A primary external cohort of single-resident homes disjoint from the development panel has since been frozen with per-home checksums and screening provenance, together with the candidate configuration and a one-shot scoring runner. That scoring has since been executed once: the median paired difference in household balanced accuracy was $+0.0091$, 95\% interval $[+0.0054, +0.0117]$, improving 37 of 43 homes, so the one component carried forward from development transfers to unfamiliar homes and unfamiliar instrumentation. Its abstention secondary outcome did not move, remaining near $10^{-3}$. The relevant question is not whether the simulator numbers reproduce exactly, but whether the qualitative findings survive: graded discrimination loss, conditional attribution value, a sensor-information frontier, non-additive sensor value, and the possibility that failure-aware weighting improves one proper score while worsening others.

\section{Reproducibility and software}

The confirmatory experiment is frozen to the provenance in
\cref{tab:provenance}, together with the pinned numerical environment recorded in
the accepted artifact.

\begin{table}[h]
\centering
\small
\caption{Frozen provenance of the accepted confirmatory run.}
\label{tab:provenance}
\begin{tabular}{@{}l p{7.2cm}@{}}
\toprule
Git revision & \texttt{bd5a9b10\allowbreak 068b69fc\allowbreak 75cc6f80\allowbreak 69044306\allowbreak 33ce7408} \\
Configuration SHA-256 & \texttt{7b48ec89\allowbreak 155f78ae\allowbreak 7965ac59\allowbreak c9b17a6b\allowbreak f164889e\allowbreak d7f38194\allowbreak c35ba157\allowbreak 6aedfae4} \\
Python & 3.11.16 \\
\bottomrule
\end{tabular}
\end{table}

The original production run generated 40 immutable shards; a later merge-only recovery reused those shards after a workflow-path bug prevented the first merge from starting. The recovered merged artifact passed the original frozen validator without rerunning any household simulation. A repository-retained result snapshot records the validator output, accepted summaries, workflow identifiers, artifact digest, and file-level hashes.

The accepted result remains simulator-derived. No field-performance claim is inferred from the validation status, and no secondary H2 sensor subset is promoted to replace the frozen primary comparison.

\section{Conclusion}

Failure-aware multimodal behavioural sensing is not validated by a single accuracy number. The frozen simulation instead reveals where the architecture succeeds and where its intuitions are incomplete. Random missingness produces graded discrimination loss, but the current abstention mechanism does not respond strongly to declining evidence. The selected five-sensor deployment is decisively not non-inferior to the ten-sensor system, although an eight-sensor secondary configuration preserves much more information. Occupancy-aware attribution helps in several contaminated scenarios but can become harmful when resident-specific evidence disappears. Failure-aware reliability weighting improves log loss while worsening balanced accuracy, Brier score, and calibration error. Sensor interactions are consistently positive for the four pre-specified pairs.

These mixed and negative findings are the main value of the confirmatory study. They narrow the claims that should be carried into real-home evaluation and prevent a simulator from being used merely to certify a preferred architecture. The next scientific boundary is therefore external validation on held-out real households, with the present negative H2 result and metric disagreements preserved rather than tuned away. That boundary is no longer entirely ahead of us. Development evidence on 22 real homes (\cref{sec:realdev}) already indicates that the simulator represents an easier problem than real instrumentation supports, and independently corroborates the H1 finding that abstention does not track evidence quality. The frozen external cohort has since been scored, and its outcome sharpens rather than softens that reading. The candidate's primary balanced-accuracy difference was positive with an interval excluding zero, so the architecture does transfer; its abstention secondary outcome stayed inert, as it did here under 40\% missingness and on the real development recordings. Three independent bodies of evidence now agree that this framework does not know when its evidence is insufficient, and the one that agrees most cheaply is this simulation, where the failure is easiest to dismiss and was seen first.

\printbibliography

\end{document}
